White-box deception probes are a leading proposal for monitoring language models, but a false answer can come from two very different sources: the model may not know the answer (hallucination), or it may know the answer and misreport it under an incentive (lie). We ask how often each happens and whether standard detectors can tell them apart. We cross knowledge (known, partial, unknown; measured by neutral sampling) with four incentive conditions (neutral, sycophantic hint, an operator-incentive prompt that never mentions lying, and an explicit instruction) on 2,975 TriviaQA questions for Llama-3.1-8B-Instruct, and score all 11,900 answers with three lie detectors and four hallucination detectors. Without pressure, 97\% of false answers come from questions the model does not fully know; under an operator incentive, 51\% are knowing misreports, and 73\% of known questions receive a false answer that the model corrects when asked for the truth in the same conversation. An Apollo-style deception probe separates lies from hallucinations with AUROC 0.89, but it reaches 0.99 from the prompt tokens alone, before any answer exists. When a knowing lie and a false answer to an unknown question share the same pressure prompt, every off-the-shelf detector is near chance (AUROC 0.50--0.60). Hallucination detectors track confidence, so they miss confident incentive lies (log-prob AUROC 0.15 for lie vs.\ hallucination). Only a composite of a falsity probe and a knowledge probe separates same-prompt lies from ignorant falsehoods in every condition (0.65--0.70). Lie-detector evaluations need same-prompt negatives that are false but uninformed, and detectors should model knowledge explicitly.
